\begin{tikzpicture}[
  x=1cm,
  y=1cm,
  font=\figurefont\footnotesize,
  source/.style={draw=FigureInput,line width=0.9pt,fill=FigureInput!9,minimum height=10mm,text width=45mm,align=left,inner sep=2mm},
  fact/.style={draw=FigureModel,line width=1pt,fill=FigureModel!10,minimum height=10mm,text width=53mm,align=center,inner sep=2mm},
  process/.style={draw=BookMuted,line width=0.8pt,fill=BookPaper,minimum height=10mm,text width=24mm,align=center,inner sep=1.5mm},
  decision/.style={draw=BookAmber,line width=1pt,fill=BookAmber!9,minimum height=12mm,text width=25mm,align=center,inner sep=1.5mm},
  result/.style={draw=FigureOutput,line width=1pt,fill=FigureOutput!7,minimum height=11mm,text width=28mm,align=center,inner sep=1.5mm},
  flow/.style={-{Stealth[length=2mm]},draw=BookInk,line width=1.1pt,rounded corners=1mm},
  evidence/.style={-{Stealth[length=1.8mm]},draw=FigureModel,line width=0.9pt,dashed,rounded corners=1mm},
  edge label/.style={fill=white,text=BookMuted,inner sep=1pt,font=\figurefont\scriptsize}
]
  \fill[white] (0,0) rectangle (16.7,6.8);
  \node[fact] (kb) at (8.3,6.1) {已有构造事实\\$(\text{林舟/E-017},\text{任职于},\text{海岬科技/E-003})$};

  \node[source,anchor=east] (sent-a) at (4.9,4.65) {句甲：海岬科技任命林舟为首席架构师。};
  \node[source,anchor=east] (sent-b) at (4.9,2.75) {句乙：林舟在论坛上批评了海岬科技的发布流程。};

  \node[process] (entity) at (6.35,3.7) {实体消歧\\确认同一实体对};
  \node[process] (mapping) at (9.25,3.7) {关系定义对应\\方向与时间};
  \node[decision] (scope) at (12.2,3.7) {当前句是否\\表达该事实？};

  \node[result] (positive) at (15.2,4.65) {句子级候选\\句甲：任职于};
  \node[result,draw=BookMuted,dashed,fill=white] (cooccur) at (15.2,2.75) {不产生句子正例\\句乙：仅共同出现};
  \node[result,text width=42mm] (bag) at (12.2,1.05) {提及组级候选\\至少一条提及可能表达关系};

  \draw[flow] (sent-a.east) -- ++(0.35,0) |- (entity.west);
  \draw[flow] (sent-b.east) -- ++(0.35,0) |- (entity.west);
  \draw[flow] (entity) -- (mapping);
  \draw[flow] (mapping) -- (scope);
  \draw[flow] (scope.east) -- ++(0.2,0) |- (positive.west);
  \draw[flow] (scope.east) -- ++(0.2,0) |- (cooccur.west);
  \draw[flow] (positive.south) -- ++(0,-0.35) -| (bag.east);

  \draw[evidence] (kb.south) -- (mapping.north);
  \draw[evidence] (kb.south east) -- ++(2.9,-0.55) |- (scope.north);

  \draw[flow] (0.35,0.65) -- (1.25,0.65);
  \node[anchor=west,text=BookMuted,font=\figurefont\scriptsize] at (1.35,0.65) {处理步骤};
  \draw[evidence] (2.95,0.65) -- (3.85,0.65);
  \node[anchor=west,text=BookMuted,font=\figurefont\scriptsize] at (3.95,0.65) {事实的证据关联};
\end{tikzpicture}
